\par
\subsection{The $N = 60$ revival and the $5 \times 5$ torus}
The longest period, $N = 60$, occurs on the $5 \times 5$,
$5 \times 10$, and $10 \times 10$ tori, and is the only period on
the pentagonal cycles. The $5 \times 5$ torus alone contributes
9 of the 182 parameter sets at
this period, making it the richest single torus by concentration.

The reflection parameters are the golden-ratio surds
$(5 - \sqrt{5})/10 \approx 0.2764$ and its complement
$(5 + \sqrt{5})/10 \approx 0.7236$. The smaller value satisfies
$\rho = 1/(\varphi + 2)$ for the golden ratio
$\varphi = (1 + \sqrt{5})/2$. The phase parameters take values in
$\{0, 2\pi/5, 4\pi/5\}$, reflecting the pentagonal symmetry of
the underlying $5$-cycle. This is the clearest instance in the
data of an inverse relationship between torus size and maximum
revival period: the smallest prime torus supports the longest
revival.
